\section{Deployment requirements}
\label{sec:deployment}

A facility deciding whether to sign asks one question about hardware: will the computer we already own
run this? The figures below are a \emph{sizing baseline} to check a machine against, not a measurement
--- planning assumptions until the prototype is benchmarked on a machine of this class. The named
processor is an illustration, not a purchase order: any part of comparable class is expected to meet
the same workload. Comparable class is settled by a benchmark, not by a spec sheet — the same test
that confirms the reference part below would confirm any candidate — so until that benchmark exists,
compare core count and vintage rather than one clock number.

\noindent\textbf{One facility node.} Each participating health center runs one. It trains on its own
records and uploads the resulting parameters; no row, identifier or label leaves the machine
(Section~\ref{sec:entities-system}).

\begin{spectab}{|>{\raggedright\arraybackslash\bfseries}p{0.145\textwidth}|L{0.315\textwidth}|L{0.44\textwidth}|}
\hline
\rowcolor{headcell}\textbf{Part} & \textbf{Planning baseline} & \textbf{Why it is set here} \\ \hline
\endfirsthead
\hline
\rowcolor{headcell}\textbf{Part} & \textbf{Planning baseline} & \textbf{Why it is set here} \\ \hline
\endhead
Processor & 6 cores, about 2.9\,GHz base; an Intel Core i5-10400 is the reference part & Training reads hundreds to low thousands of rows against a model specified to be simple on purpose (Section~\ref{sec:scope}). A model that needed a workstation would undercut that position. \\ \hline
\rowcolor{rowshade}Memory & 8\,GB & One operating system, one interpreter and one model. The model persists between rounds; nothing else is held. The operating system is the binding case, not the model. \\ \hline
Storage & 256\,GB, at least 40\,GB free & Facility records stay on this machine. Their retention is settled by the data-sharing agreement and C3 (Section~\ref{sec:conformance}), not by this section. \\ \hline
\rowcolor{rowshade}Graphics & Integrated only & No step in this system is specified to need a discrete graphics processor. A machine that has one is not disqualified. \\ \hline
Operating system & Windows 11 Pro 64-bit, or Ubuntu 22.04 LTS & Both are intended targets and both must be exercised before either is called supported. Windows~11 is limited to Microsoft's published supported-processor lists (Intel 8th generation and newer, AMD Ryzen 2000 and newer), which is where the floor comes from; the compatibility check is the authority, not this table. \\ \hline
\rowcolor{rowshade}Network & 512\,kbps upload, with interruptions & A round sends one model file: about sixteen seconds at 512\,kbps for a one-megabyte file, before protocol overhead. Section~\ref{sec:metrics} prices a whole round; it does not establish that the link dominates one. \\ \hline
Power & 600\,VA uninterruptible supply & An interrupted write can damage records the facility holds. Whether the supply is the facility's cost or the study's is not settled here. \\ \hline
\end{spectab}

\noindent\textbf{The aggregator.} One machine for the whole pilot, not one per health center. It receives
parameter updates and keeps the model versions and the run record; it is the same order of computer as
a node, not a server.

\begin{spectab}{|>{\raggedright\arraybackslash\bfseries}p{0.145\textwidth}|L{0.315\textwidth}|L{0.44\textwidth}|}
\hline
\rowcolor{headcell}\textbf{Part} & \textbf{Planning baseline} & \textbf{Why it is set here} \\ \hline
\endfirsthead
\hline
\rowcolor{headcell}\textbf{Part} & \textbf{Planning baseline} & \textbf{Why it is set here} \\ \hline
\endhead
Processor & 4 cores or more; the same reference part, or a 4-vCPU virtual machine & It combines the updates it receives and writes them to disk. A hosted virtual machine is an expected arrangement, not a compromise. \\ \hline
\rowcolor{rowshade}Memory & 16\,GB & The global model and every facility's update are held at once, and that number grows with the facility count $N$ (Section~\ref{sec:parameters}). The 16\,GB figure assumes at most ten facilities and a model of a few megabytes. \\ \hline
Storage & 512\,GB, at least 250\,GB free & One model file per round, so versions cost about $R \times$ model size: at 200 rounds and 1\,MB, roughly 200\,MB. That is an illustration, not a measured setting, and it assumes a model of a few megabytes over a round budget of a few hundred. \\ \hline
\rowcolor{rowshade}Operating system & Ubuntu 22.04 LTS & It is intended to run unattended for the length of a study. Windows Server is untested and not claimed. \\ \hline
Network & Wired, 5\,Mbps or more, reachable by every facility & Every upload arrives here, so it has to be the dependable end of an unreliable network. \\ \hline
\rowcolor{rowshade}Backup & A copy off the machine at least nightly & The model versions are the study's evidence; losing them loses the run. The cadence is a planning figure. \\ \hline
\end{spectab}

\noindent\textbf{Checking a machine already on the desk.} On Windows, the processor and installed memory
are under \textsf{Settings} $\rightarrow$ \textsf{System} $\rightarrow$ \textsf{About}, and free space
under \textsf{This PC}; Windows~11 eligibility is answered by Microsoft's own compatibility check. For
the link, the upload rate matters, not the download figure a speed test quotes first. A machine that
meets every row is expected to need nothing bought, and one that misses a row is expected to need
that row fixed rather than a new computer; the benchmark is what confirms either reading.


\noindent\textbf{Data requested from the City Health Office.} The study needs one kind of record:
what left the store. Each entry should carry five facts, and the
office's own file already carries most of them under whatever names it uses.

\begin{enumerate}[leftmargin=1.6em,itemsep=0.25em]
  \item \textbf{Facility.} The office's way of telling one health center from another: a name, a
        code, or both. National facility and barangay codes if they are held; whether they are is
        OD-6.
  \item \textbf{Medicine.} Generic name, dosage form, strength and unit, as the facility stocks it
        (Section~\ref{sec:entities-domain}). An internal product code or a brand name may travel
        alongside. Where the office records only codes or brand names, converting them to generic
        names can be done by the office or by the study, whichever suits the office --- the request
        is that each medicine carry a generic name by the time the file is used.
  \item \textbf{Period.} The span each entry covers, stated once for the file. A day, a week and a
        month all work; which one is used remains OD-2, and all this asks is that the file say which.
  \item \textbf{Quantity.} How much left the store in that period, and the unit it is counted in
        (tablets, bottles, boxes). A closing stock balance instead of a dispensed quantity can also
        be worked with.
  \item \textbf{Completeness.} A marker on entries that were estimated or reconstructed, or a note
        saying the office keeps none. A gap that is marked is usable data; a gap quietly filled in
        is not.
\end{enumerate}

\noindent
Column names, codes and file layout stay the office's. A system export, a spreadsheet or a written
log typed up afterwards all serve. If the facts sit under unfamiliar headings, or across two files
that join on facility and date, that is acceptable: what would not be acceptable is a file from
which those five facts cannot be recovered.

\noindent
History is whatever continuous run the office holds, with its start and end stated. A short run is
accepted and recorded as short; the study does not ask for a period the records cannot cover, and a
longer run is better than a short one because demand has seasons. Patient identities, prescriptions,
staff names and prices are not requested, and a file containing them will be reduced to the five
facts above before it is used (C1).

\noindent
The mapped file becomes one dataset version with a stated purpose and date (C4). Until an agreement
exists and the real file is provided, every figure in this annex comes from a synthetic dataset built
to this shape, and is labelled synthetic (C5).

\noindent\textbf{What this section does not settle.} Who hosts the aggregator, and therefore who pays
for it and for its backup, is open decision OD-5; the shape of the data request follows OD-2 and
OD-6 (both Section~\ref{sec:open}). Whether power protection is
the facility's obligation or the study's is undecided. And how the quality metric of
Section~\ref{sec:metrics} is computed across facilities in a real deployment is open decision OD-8:
the specification describes the computation as the study runs it, on held-out rows pooled inside one
process. A deployment would instead need some exchange --- a score per facility, or a sufficient statistic
--- and neither is a row or a parameter: Section~\ref{sec:scope} promises that only parameters
cross, so that promise is narrower than deployment evaluation would need. Answer OD-8 before this section's figures are
treated as sufficient.
